\section{Projection formula}
\label{section-projection-formula}

\noindent
In this section we collect variants of the projection formula.
The most basic version is Lemma \ref{lemma-projection-formula}.
After we state and prove it, we discuss a more general version
involving perfect complexes.

\begin{lemma}
\label{lemma-injective-tensor-finite-locally-free}
Let $X$ be a ringed space.
Let $\mathcal{I}$ be an injective $\mathcal{O}_X$-module.
Let $\mathcal{E}$ be an $\mathcal{O}_X$-module.
Assume $\mathcal{E}$ is finite locally free on $X$, see
Modules, Definition \ref{modules-definition-locally-free}.
Then $\mathcal{E} \otimes_{\mathcal{O}_X} \mathcal{I}$ is
an injective $\mathcal{O}_X$-module.
\end{lemma}

\begin{proof}
This is true because under the assumptions of the lemma we have
$$
\Hom_{\mathcal{O}_X}(\mathcal{F},
\mathcal{E} \otimes_{\mathcal{O}_X} \mathcal{I})
=
\Hom_{\mathcal{O}_X}(
\mathcal{F} \otimes_{\mathcal{O}_X} \mathcal{E}^\vee, \mathcal{I})
$$
where
$\mathcal{E}^\vee = \SheafHom_{\mathcal{O}_X}(\mathcal{E}, \mathcal{O}_X)$
is the dual of $\mathcal{E}$ which is finite locally free also. Since tensoring
with a finite locally free sheaf is an exact functor we win by
Homology, Lemma \ref{homology-lemma-characterize-injectives}.
\end{proof}

\begin{lemma}
\label{lemma-projection-formula}
Let $f : X \to Y$ be a morphism of ringed spaces.
Let $\mathcal{F}$ be an $\mathcal{O}_X$-module.
Let $\mathcal{E}$ be an $\mathcal{O}_Y$-module.
Assume $\mathcal{E}$ is finite locally free on $Y$, see
Modules, Definition \ref{modules-definition-locally-free}.
Then there exist isomorphisms
$$
\mathcal{E} \otimes_{\mathcal{O}_Y} R^qf_*\mathcal{F}
\longrightarrow
R^qf_*(f^*\mathcal{E} \otimes_{\mathcal{O}_X} \mathcal{F})
$$
for all $q \geq 0$. In fact there exists an isomorphism
$$
\mathcal{E} \otimes_{\mathcal{O}_Y} Rf_*\mathcal{F}
\longrightarrow
Rf_*(f^*\mathcal{E} \otimes_{\mathcal{O}_X} \mathcal{F})
$$
in $D^{+}(Y)$ functorial in $\mathcal{F}$.
\end{lemma}

\begin{proof}
Choose an injective resolution $\mathcal{F} \to \mathcal{I}^\bullet$
on $X$. Note that $f^*\mathcal{E}$ is finite locally free also, hence
we get a resolution
$$
f^*\mathcal{E} \otimes_{\mathcal{O}_X} \mathcal{F}
\longrightarrow
f^*\mathcal{E} \otimes_{\mathcal{O}_X} \mathcal{I}^\bullet
$$
which is an injective resolution by
Lemma \ref{lemma-injective-tensor-finite-locally-free}.
Apply $f_*$ to see that
$$
Rf_*(f^*\mathcal{E} \otimes_{\mathcal{O}_X} \mathcal{F})
=
f_*(f^*\mathcal{E} \otimes_{\mathcal{O}_X} \mathcal{I}^\bullet).
$$
Hence the lemma follows if we can show that
$f_*(f^*\mathcal{E} \otimes_{\mathcal{O}_X} \mathcal{F}) =
\mathcal{E} \otimes_{\mathcal{O}_Y} f_*(\mathcal{F})$ functorially
in the $\mathcal{O}_X$-module $\mathcal{F}$. This is clear when
$\mathcal{E} = \mathcal{O}_Y^{\oplus n}$, and follows in general
by working locally on $Y$. Details omitted.
\end{proof}

\noindent
Let $f : X \to Y$ be a morphism of ringed spaces. 
Let $E \in D(\mathcal{O}_X)$ and $K \in D(\mathcal{O}_Y)$.
Without any further assumptions there is a map
\begin{equation}
\label{equation-projection-formula-map}
Rf_*E \otimes^\mathbf{L}_{\mathcal{O}_Y} K
\longrightarrow
Rf_*(E \otimes^\mathbf{L}_{\mathcal{O}_X} Lf^*K)
\end{equation}
Namely, it is the adjoint to the canonical map
$$
Lf^*(Rf_*E \otimes^\mathbf{L}_{\mathcal{O}_Y} K) =
Lf^*Rf_*E \otimes^\mathbf{L}_{\mathcal{O}_X} Lf^*K
\longrightarrow
E \otimes^\mathbf{L}_{\mathcal{O}_X} Lf^*K
$$
coming from the map $Lf^*Rf_*E \to E$ and Lemmas
\ref{lemma-pullback-tensor-product} and \ref{lemma-adjoint}.
A reasonably general version of the projection formula is the following.

\begin{lemma}
\label{lemma-projection-formula-perfect}
Let $f : X \to Y$ be a morphism of ringed spaces.
Let $E \in D(\mathcal{O}_X)$ and $K \in D(\mathcal{O}_Y)$.
If $K$ is perfect, then
$$
Rf_*E \otimes^\mathbf{L}_{\mathcal{O}_Y} K =
Rf_*(E \otimes^\mathbf{L}_{\mathcal{O}_X} Lf^*K)
$$
in $D(\mathcal{O}_Y)$.
\end{lemma}

\begin{proof}
To check (\ref{equation-projection-formula-map}) is an isomorphism
we may work locally on $Y$, i.e., we have to find an open covering
$\{V_j \to Y\}$
such that the map restricts to an isomorphism on $V_j$. By definition
of perfect objects, this means we may assume $K$ is represented by
a strictly perfect complex of $\mathcal{O}_Y$-modules.
Note that, completely generally, the statement is true for
$K = K_1 \oplus K_2$, if and only if the statement is true for
$K_1$ and $K_2$. Hence we may assume $K$ is a finite
complex of finite free $\mathcal{O}_Y$-modules.
In this case a simple argument involving stupid truncations reduces
the statement to the case where $K$ is represented by a finite
free $\mathcal{O}_Y$-module. Since the statement is invariant
under finite direct summands in the $K$ variable, we conclude
it suffices to prove it for $K = \mathcal{O}_Y[n]$
in which case it is trivial.
\end{proof}

\noindent
Here is a case where the projection formula is true in complete
generality.

\begin{lemma}
\label{lemma-projection-formula-closed-immersion}
Let $f : X \to Y$ be a morphism of ringed spaces such that $f$ is a
homeomorphism onto a closed subset. Then
(\ref{equation-projection-formula-map}) is an isomorphism always.
\end{lemma}

\begin{proof}
Since $f$ is a homeomorphism onto a closed subset, the functor $f_*$
is exact (Modules, Lemma \ref{modules-lemma-i-star-exact}). Hence
$Rf_*$ is computed by applying $f_*$ to any representative complex.
Choose a K-flat complex $\mathcal{K}^\bullet$ of $\mathcal{O}_Y$-modules
representing $K$ and choose any complex $\mathcal{E}^\bullet$
of $\mathcal{O}_X$-modules representing $E$. Then
$Lf^*K$ is represented by $f^*\mathcal{K}^\bullet$ which is
a K-flat complex of $\mathcal{O}_X$-modules
(Lemma \ref{lemma-pullback-K-flat}). Thus the right hand side of
(\ref{equation-projection-formula-map}) is represented by
$$
f_*\text{Tot}(\mathcal{E}^\bullet
\otimes_{\mathcal{O}_X} f^*\mathcal{K}^\bullet)
$$
By the same reasoning we see that the left hand side is represented by
$$
\text{Tot}(f_*\mathcal{E}^\bullet \otimes_{\mathcal{O}_Y} \mathcal{K}^\bullet)
$$
Since $f_*$ commutes with direct sums
(Modules, Lemma \ref{modules-lemma-i-star-right-adjoint})
it suffices to show that
$$
f_*(\mathcal{E} \otimes_{\mathcal{O}_X} f^*\mathcal{K}) =
f_*\mathcal{E} \otimes_{\mathcal{O}_Y} \mathcal{K}
$$
for any $\mathcal{O}_X$-module $\mathcal{E}$ and $\mathcal{O}_Y$-module
$\mathcal{K}$. We will check this by checking on stalks.
Let $y \in Y$. If $y \not \in f(X)$, then the stalks
of both sides are zero. If $y = f(x)$, then we see that we have to show
$$
\mathcal{E}_x \otimes_{\mathcal{O}_{X, x}}
(\mathcal{O}_{X, x} \otimes_{\mathcal{O}_{Y, y}} \mathcal{F}_y) =
\mathcal{E}_x \otimes_{\mathcal{O}_{Y, y}} \mathcal{F}_y
$$
(using Sheaves, Lemma \ref{sheaves-lemma-stalks-closed-pushforward}
and Lemma \ref{sheaves-lemma-stalk-pullback-modules}).
This equality holds and therefore the lemma has been proved.
\end{proof}

\begin{remark}
\label{remark-compatible-with-diagram}
The map (\ref{equation-projection-formula-map}) is compatible with the
base change map of Remark \ref{remark-base-change} in the following sense.
Namely, suppose that
$$
\xymatrix{
X' \ar[r]_{g'} \ar[d]_{f'} &
X \ar[d]^f \\
Y' \ar[r]^g &
Y
}
$$
is a commutative diagram of ringed spaces. 
Let $E \in D(\mathcal{O}_X)$ and $K \in D(\mathcal{O}_Y)$.
Then the diagram
$$
\xymatrix{
Lg^*(Rf_*E \otimes^\mathbf{L}_{\mathcal{O}_Y} K) \ar[r]_p \ar[d]_t &
Lg^*Rf_*(E \otimes^\mathbf{L}_{\mathcal{O}_X} Lf^*K) \ar[d]_b \\
Lg^*Rf_*E \otimes^\mathbf{L}_{\mathcal{O}_{Y'}} Lg^*K \ar[d]_b &
Rf'_*L(g')^*(E \otimes^\mathbf{L}_{\mathcal{O}_X} Lf^*K) \ar[d]_t \\
Rf'_*L(g')^*E \otimes^\mathbf{L}_{\mathcal{O}_{Y'}} Lg^*K \ar[rd]_p &
Rf'_*(L(g')^*E \otimes^\mathbf{L}_{\mathcal{O}_{Y'}} L(g')^*Lf^*K) \ar[d]_c \\
& Rf'_*(L(g')^*E \otimes^\mathbf{L}_{\mathcal{O}_{Y'}} L(f')^*Lg^*K)
}
$$
is commutative. Here arrows labeled $t$ are gotten by an application of
Lemma \ref{lemma-pullback-tensor-product}, arrows labeled $b$ by an
application of Remark \ref{remark-base-change}, arrows labeled $p$
by an application of (\ref{equation-projection-formula-map}), and
$c$ comes from $L(g')^* \circ Lf^* = L(f')^* \circ Lg^*$.
We omit the verification.
\end{remark}















